\begin{textsample}{POS Dim 5 – human – Score 1.00 – mg/mg\_ef\_ciencias\_humanas\_intro\_2\_p1.txt}  \label{ex:f5_pos_001}
Minas Gerais é um Estado de encontro de povos e, por isso, apresenta ampla e diversa expressão cultural em seus mais variados registros, tanto no passado, quanto no presente.
O reconhecimento de nossa diversidade e o respeito às expressões culturais deve fazer parte do cotidiano da sala de aula de forma a contribuir para a convivência social autônoma e respeitosa.
As muitas \textbf{Minas} Gerais já evocadas pelo grande Guimarães Rosa são palco de um jeitinho próprio de cozinhar, de cultivar a terra, de conversar, de dançar, de expressar, de ser, que denotam características próprias de um povo que em sua diversidade consegue construir sua identidade.
A identidade mineira é representada pela grande diversidade sociocultural e natural, distribuída nos oitocentos e cinquenta e três municípios que compõem o Estado, no qual convivem diferentes grupos como os sertanejos, as populações ribeirinhas, povos indígenas, remanescentes de quilombos, povos ciganos, entre outros.
Cada localidade tem sua própria dinâmica, o que inclui o patrimônio cultural material e imaterial, representado por sítios arqueológicos, arquitetura barroca, tradições folclóricas, gastronomia, diversidade religiosa, entre outros.

% matched lemmas: minas
\end{textsample}
